\section{Measure and integration from coherent finite observations}
\label{sec:measure}

The repository already has an effective expectation theory, a dominated
convergence statement, and a Vitali statement
\cite{repo-expectation-packet}.  Its current object-language development is
substantial: \texttt{FiniteKernel.pr\_evaluator} makes cross-multiplied
expectation comparisons quantifier-free; \texttt{TowerCommonDen.regroup}
proves the numerator regrouping used by the tower property; and
\texttt{TruncationEstimate.truncation} proves a finite inequality with a
common dominator
\cite{repo-finite-kernel,repo-tower-common-den,repo-truncation}.
The new work here concerns the passage from these finite identities to
coherent measure-theoretic objects.  It supplies a stable convergence
observable, constructs an integrable limit rather than presupposing one,
extracts amplitude cutoffs from an existing integrable name, and extends
finite conditional averages with an explicit approximation rate.

There is also a specific local correction.  The earlier packet's human
Vitali proof uses
\[
 |x-y|\le |x-y|\mathbf1_{|x|\le K,\ |y|\le K}
       +|x|\mathbf1_{|x|>K}+|y|\mathbf1_{|y|>K}.
\]
For $K>0$, the values $x=2K$, $y=-K$ give $3K\le2K$, so that displayed
inequality is false.  This does not affect the different, common-dominator
inequality in the checked \texttt{TruncationEstimate} module.  The clipping
estimate proved below repairs the Vitali argument and also avoids computing
hard threshold events from approximations to real values.

Throughout this section put $\delta_k=2^{-k}$.  Every input evaluator,
refinement map and modulus mentioned as effective is a \emph{fixed}
primitive-recursive program. For an iPRA export, its displayed numerical
contracts must be supplied as object-theory derivations or retained as
explicit premises. Numerical problem parameters may vary. This is a
schema about fixed programs, not a claim that arbitrary codes for total
programs have a primitive-recursive universal evaluator.

\subsection{A common probability algebra and its integrable completion}

\begin{definition}[Coherent rational probability presentation]
For each $n\in\N$, let $A_n=\{0,\ldots,a(n)-1\}$ be a nonempty finite set
of atoms, with rational weights $w_{n,i}\ge0$ summing to one.  For $m\le n$
there is a map $\pi_{m,n}:A_n\to A_m$, satisfying
\begin{align}
 \pi_{n,n}(i)&=i,&
 \pi_{\ell,n}&=\pi_{\ell,m}\circ\pi_{m,n}\quad(\ell\le m\le n),
 \label{eq:m-refinement}\\
 w_{m,i}&=\sum_{\substack{j<a(n)\\\pi_{m,n}(j)=i}}w_{n,j}.
 \label{eq:m-mass}
\end{align}
The cardinalities, weights and maps have fixed PR evaluators.  Their range,
normalization and compatibility equations are part of the presentation.
\end{definition}

A simple function is a code $s=(n,(s_i)_{i<a(n)})$, with rational values.
To compare $s$ at level $n$ with $t$ at level $m$, lift both lists to
$r=\max(m,n)$ using $\pi_{n,r}$ and $\pi_{m,r}$.  On that common level set
\begin{equation}
 I(s)=\sum_iw_i s_i,\qquad
 \|s\|_1=\sum_iw_i|s_i|,\qquad
 d_1(s,t)=\sum_iw_i|s_i-t_i|.
 \label{eq:m-simple}
\end{equation}
Equation~\eqref{eq:m-mass} proves that lifting again changes none of these
values.  It also proves independence of the chosen common refinement.
Finite addition, scalar multiplication, minimum, maximum and absolute value
are computed on a common level.  Zero-weight cells are retained in codes
but have zero distance.  Thus this is a rational vector lattice with a
pseudometric, rather than a demand for decidable equality of functions.

The coherent tower matters.  Marginal probability laws alone do not
determine a common $L^1$ geometry: two fair signs can be equal or opposite,
giving respectively zero or two as the expectation of their absolute
difference.  The maps \eqref{eq:m-refinement} choose the common probability
algebra explicitly and keep every operation on that algebra.

This presentation has concrete continuous models.  For dyadic intervals
in $[0,1)$, take $a(n)=2^n$, $w_{n,i}=2^{-n}$ and
$\pi_{m,n}(i)=\lfloor i/2^{n-m}\rfloor$.  More generally, every tower above
has an explicit interval realization.  Give each atom of $A_0$ a half-open
interval of its stated length.  Inside the interval for $i\in A_n$, put
the children $j$ with $\pi_{n,n+1}(j)=i$ in increasing order, with lengths
$w_{n+1,j}$.  The left endpoint of child $j$ is the parent's left endpoint
plus the sum of the weights of its earlier siblings.  Compatibility makes
these children fill the parent interval exactly.  Zero weights give empty
intervals.  This recursive rational construction realizes all finite
integrals as ordinary step-function integrals.  No infinite branch is
selected.  In a usual semantic $L^1$ space the completion below is precisely
the closure of these step functions; asserting density in a larger space
is a separate adequacy statement.  The interval realization alone neither
codes all measurable subsets nor supplies a rate for arbitrary countable
unions.

\begin{definition}[Regular integrable name]
A regular $L^1$ name is a fixed PR sequence of simple-function codes
$s_k$ satisfying
\begin{equation}
 d_1(s_k,s_\ell)\le\delta_k+\delta_\ell
 \qquad(k,\ell\in\N).
 \label{eq:m-regular}
\end{equation}
Its semantic value $f$ is the corresponding completion element; in
particular $\|f-s_k\|_1\le\delta_k$.  Equivalent names denote the same
completion element.  Internally one works with
\eqref{eq:m-regular} and finite comparisons of the actual codes, not with
a new sort of arbitrary measurable functions.
\end{definition}

For example, the distance between two represented elements is itself a
regular scalar name: evaluate $d_1(s_{k+1},t_{k+1})$; the shift absorbs
both input errors.  Addition similarly uses $s_{k+1}+t_{k+1}$ at output
precision $k$.  Equality means that every requested
distance observation can be made arbitrarily small.  It is an arithmetical
formula, not an algorithm which tests exact real equality.

\begin{lemma}[Stable extension of finite observables]
\label{lem:m-extension}
The integral and norm extend to integrable names, with
\begin{equation}
 |I(f)-I(g)|\le\|f-g\|_1,\qquad
 \bigl|\|f\|_1-\|g\|_1\bigr|\le\|f-g\|_1.
 \label{eq:m-integral-stable}
\end{equation}
Any finite operation $T$ satisfying
$d_1(Ts,Tt)\le d_1(s,t)$ extends by the same precision schedule.  This
applies to $s\mapsto |s|$, $s\mapsto s^+=\max(s,0)$, and clipping to a
fixed rational interval.  Finite identities between such operations
remain true on the completion.
\end{lemma}

\begin{proof}
The finite triangle inequality gives
$|\sum_iw_i(s_i-t_i)|\le\sum_iw_i|s_i-t_i|$.  The reverse triangle
inequality gives the second assertion in \eqref{eq:m-integral-stable}.
Consequently the finite values on a regular name form regular scalar
names.  A nonexpansive $T$ carries \eqref{eq:m-regular} to the identical
inequality for $Ts_k$.  If two input names represent the same element,
their image names do also, by this inequality.  A finite identity is an
identity at every common approximation level; each side's error then tends
to zero with its displayed schedule.  This proves existence,
well-definedness and preservation, without choosing representatives at
individual sample points.
\end{proof}

The completed order is $f\le g$ when $\|(f-g)^+\|_1=0$.  It is stable
under limits by the same Lipschitz argument.  Multiplication requires its
own contract: a bounded multiplier $|h|\le M$ satisfies
$\|hf-hg\|_1\le M\|f-g\|_1$.  Products of two unrestricted $L^1$
elements are not part of this interface.

Finite events are subsets of a level's atom set.  Their metric is
$d_\mu(B,C)=\mu(B\mathbin\triangle C)$; their indicator functions embed
isometrically in the simple $L^1$ space.  Complement preserves this metric,
and union and intersection change by at most the sum of the input errors.
Thus regular event names form an effective measure algebra.  For a
uniformly PR-presented sequence of event names, a countable union can be
constructed when its finite partial unions $B^{<n}$ have a PR continuity
modulus
\[
 m\ge n\ge R(k)\quad\Longrightarrow\quad
 \mu(B^{<m}\setminus B^{<n})\le\delta_k.
\]
Replace $R$ by its bounded monotone envelope if needed.  To make the
diagonal explicit, take $N=R(k+2)$, approximate each of the
first $N$ events to error
$2^{-k-2-\lceil\log_2(N+1)\rceil}$, and take their finite union.  Its
total sampling error is at most $\delta_{k+2}$.  Together with the
displayed continuity bound this gives a regular output name by the same
three-error estimate as \eqref{eq:m-output-name} below.  All logarithms
here mean bounded integer binary-length calculations.  For disjoint
events this premise is exactly a supplied tail estimate for their mass
series; finite additivity and continuity of $\mu$ give countable
additivity for this represented union.  An arbitrary sequence of event
codes does not itself supply the displayed modulus.

\subsection{Stable observations of convergence in measure}

For $K\ge0$, define the scalar clipping operation and excess by
\[
 c_K(x)=\max(-K,\min(K,x)),\qquad
 p_K(x)=(|x|-K)^+.
\]
For integrable elements set
\begin{equation}
 T_K(f)=I(p_K(f))=\|f-c_K(f)\|_1,
 \qquad
 \rho(f,g)=I\bigl(\min(1,|f-g|)\bigr).
 \label{eq:m-stable-observers}
\end{equation}
Both expressions have evaluators inherited from finite rational lists.
The function $p_K$ is 1-Lipschitz, as is $x\mapsto\min(1,|x|)$.
Consequently
\begin{equation}
 |T_K(f)-T_K(g)|\le\|f-g\|_1,\quad
 |\rho(f,g)-\rho(f',g')|
 \le\|f-f'\|_1+\|g-g'\|_1.
 \label{eq:m-observer-lipschitz}
\end{equation}
These bounds are the reason finite approximation errors can be transported
through the observers.

\begin{lemma}[Clipping estimate]
\label{lem:m-clipping}
For $K\ge1$ and integrable $f,g$,
\begin{equation}
 \|f-g\|_1\le T_K(f)+T_K(g)+2K\rho(f,g).
 \label{eq:m-clipping}
\end{equation}
Also $\rho(f,g)\le\|f-g\|_1$, and $\rho$ satisfies the triangle inequality.
\end{lemma}

\begin{proof}
For rational $x,y$, the clipping map is nondecreasing and 1-Lipschitz,
has range $[-K,K]$, and $|x-c_K(x)|=p_K(x)$.  Hence
\begin{align*}
 |x-y|&\le p_K(x)+p_K(y)+|c_K(x)-c_K(y)|,\\
 |c_K(x)-c_K(y)|&\le\min(|x-y|,2K)
                         \le2K\min(1,|x-y|).
\end{align*}
For the last inequality, split the rational cases $|x-y|\le1$ and
$|x-y|\ge1$.  Multiply by the nonnegative weights and add.  This proves
\eqref{eq:m-clipping} on every finite common refinement.  For names,
replace each argument by a simple approximation and use
\eqref{eq:m-observer-lipschitz}; the total error is at most a fixed rational
multiple of the two approximation errors, so it has an explicit PR
precision schedule.  This establishes the same inequality in the
completion.  Finally,
$\min(1,a+b)\le\min(1,a)+\min(1,b)$ for $a,b\ge0$, which gives the
triangle inequality for $\rho$ after the scalar triangle inequality.
\end{proof}

The metric $\rho$ expresses convergence in measure without a threshold
decision.  In an ordinary probability-space interpretation, for
$0<\eta\le1$ the finite inequalities give
\begin{equation}
 \eta\,\mu(|f-g|>\eta)\le\rho(f,g)
 \le\eta+\mu(|f-g|>\eta).
 \label{eq:m-probability-comparison}
\end{equation}
Thus a two-parameter convergence-in-probability rate converts to a
$\rho$-rate, and conversely.  In the abstract numerical presentation, the
hard event in \eqref{eq:m-probability-comparison} is used only when its
measure has separately been justified.  The algorithm itself evaluates
the continuous truncated distance in \eqref{eq:m-stable-observers}.
An $L^1$ approximation alone does not compute the measure of a sharp
threshold event: constants just below and above a threshold can be
arbitrarily close in $L^1$ while their threshold-event probabilities are
zero and one.

One can also complete the simple space using $\rho$.  Call that completion
$L^0$ for this presentation.  Since $\rho\le d_1$, there is a canonical
map $L^1\to L^0$.  It is injective.  Indeed, for any fixed integrable $h$,
choose a simple $s$ with $\|h-s\|_1\le\varepsilon$ and
$K\ge\max(1,\|s\|_\infty)$.  Then $T_K(h)\le\varepsilon$ by
\eqref{eq:m-observer-lipschitz}.  If $\rho(h,0)=0$, the clipping estimate
gives $\|h\|_1\le\varepsilon$ for every supplied precision.  The theorem
below explicitly constructs an $L^1$ preimage of an $L^0$ Cauchy limit.

\subsection{A quantitative Vitali theorem that constructs its limit}

\begin{theorem}[Uniform-integrability lifting]
\label{thm:m-vitali}
Let $f_n$ be a uniformly presented family of integrable names: a fixed PR
program $S(n,j)$ returns a simple code with
$\|S(n,j)-f_n\|_1\le\delta_j$.  Suppose that monotone PR functions
$U,R:\N\to\N$, with $U(k)\ge1$, satisfy
\begin{align}
 T_{U(k)}(f_n)&\le\delta_k &&\text{for every }n,k,
 \label{eq:m-ui}\\
 \rho(f_m,f_n)&\le\delta_k &&\text{if }m,n\ge R(k).
 \label{eq:m-rho-cauchy}
\end{align}
For an integer $b\ge1$ put
\[
 \ell(b)=\min\{j\le b: b\le2^j\}.
\]
Define the explicit precision program
\begin{equation}
 q(k)=k+2+\ell\bigl(2U(k+2)\bigr),\qquad
 N(k)=R(q(k)).
 \label{eq:m-vitali-rate}
\end{equation}
Then there is a regular integrable name $f$, produced by a PR program,
such that
\begin{equation}
 n\ge N(k)\quad\Longrightarrow\quad
 \|f_n-f\|_1\le\tfrac34\delta_k.
 \label{eq:m-vitali-convergence}
\end{equation}
In particular the integrals converge with this rate.  The limit satisfies
$T_{U(k)}(f)\le\delta_k$; no integrable limit or tail bound for a limit is
an input assumption.
\end{theorem}

\begin{proof}
The bounded definition of $\ell$ is total because $2^b\ge b$; this finite
inequality is proved by quantifier-free induction.  Consequently
$2U(k+2)2^{-\ell(2U(k+2))}\le1$.  If $m,n\ge N(k)$, apply
Lemma~\ref{lem:m-clipping} with $K=U(k+2)$ and use the two input contracts:
\begin{align}
 \|f_m-f_n\|_1
 &\le2\delta_{k+2}+2U(k+2)\delta_{q(k)}
 \le3\delta_{k+2}=\tfrac34\delta_k.
 \label{eq:m-pair-cauchy}
\end{align}
This derives an $L^1$ Cauchy modulus before introducing any limit.
Monotonicity of $U,R,\ell$ makes $N$ monotone.  Nonmonotone supplied
moduli can first be replaced by their bounded maxima over initial
segments; both contracts remain valid.

Now define an actual simple-code output program by
\begin{equation}
 b_k=S(N(k+2),k+2).
 \label{eq:m-output-name}
\end{equation}
If $j\ge k$, both family indices in $b_k,b_j$ are at least $N(k+2)$.
Their represented values have distance at most $\delta_{k+2}$ by
\eqref{eq:m-pair-cauchy}.  Adding the two simple-approximation errors gives
\[
 d_1(b_k,b_j)\le2\delta_{k+2}+\delta_{j+2}
              \le\tfrac34\delta_k\le\delta_k+\delta_j.
\]
Thus $b_k$ is a regular name, and defines $f$.  For $n\ge N(k)$ and
all sufficiently large $j$, \eqref{eq:m-pair-cauchy} gives
\[
 \|f_n-b_j\|_1\le\tfrac34\delta_k+\delta_{j+2}.
\]
The definition of completion proves \eqref{eq:m-vitali-convergence}.
Equivalently, each rational observation is obtained by choosing $j$ large
enough for its specified remaining error; there is no search for a
convergence stage.

For each fixed $k$, continuity of the tail observable yields
\[
 T_{U(k)}(f)\le T_{U(k)}(f_n)+\|f-f_n\|_1
             \le\delta_k+\|f-f_n\|_1.
\]
Substitute $n=N(j)$ and increase the requested numerical precision $j$.
This proves $T_{U(k)}(f)\le\delta_k$.  Only the explicit Lipschitz bound
for $T_K$ is used; neither Fatou's lemma nor a separately assumed
uniform-integrability property of an unknown limit is needed.
Finally \eqref{eq:m-integral-stable} gives the same estimate for
$|I(f_n)-I(f)|$.
\end{proof}

This is a quantitative inverse to the inclusion $L^1\to L^0$ on a family
equipped with $U$.  In particular, if an independently presented
$L^0$ element is the $\rho$-limit of $f_n$, it equals the image of the
constructed integrable $f$, by the triangle inequality.  The conclusion
does not assume that the independently presented limit was already
integrable.

\begin{proposition}[A named dominator supplies its own cutoff]
\label{prop:m-domination}
Suppose $g\ge0$ has a regular integrable name $g_j$, and $|f_n|\le g$ for
all $n$ in the completed lattice order.  Define
\begin{equation}
 U(k)=\max\left(1,
       \max_{j\le k}\left\lceil\|g_j\|_\infty\right\rceil\right).
 \label{eq:m-derived-ui}
\end{equation}
Then \eqref{eq:m-ui} holds.  Thus, together with the $\rho$ Cauchy rate,
these hypotheses suffice for Theorem~\ref{thm:m-vitali}; a separate
amplitude-tail oracle for $g$ is unnecessary.
\end{proposition}

\begin{proof}
For $K\ge\|g_k\|_\infty$, $p_K(g_k)=0$.  By
\eqref{eq:m-observer-lipschitz},
$T_K(g)\le\|g-g_k\|_1\le\delta_k$.
The scalar inequality
\[
 (|x|-K)^+\le (y-K)^+ + (|x|-y)^+
\]
shows, after integration and stable passage to names, that $|f_n|\le g$
implies $T_K(f_n)\le T_K(g)$.  Apply this with $K=U(k)$.
The maximum and ceiling in \eqref{eq:m-derived-ui} inspect finite rational
tables and are PR.  They do not compute an essential supremum of $g$.
The constructed limit remains dominated: the norm of
$(|f|-g)^+$ is at most $\|f-f_n\|_1$ and hence is zero.
\end{proof}

The excess convention for $T_K$ is equivalent, with explicit constants,
to the usual uniform-integrability convention when sharp tail integrals
are available.  Pointwise,
\[
 (|x|-K)^+\le |x|\mathbf1_{|x|>K},\qquad
 |x|\mathbf1_{|x|>2K}\le2(|x|-K)^+.
\]
The excess is preferable as the primitive observer because it is
Lipschitz.  General quantitative and metastable convergence work provides
broader context for making these data explicit
\cite{neri-pischke,avigad-dean-rute}; the displayed schedules and proofs here
follow directly from finite rational inequalities.

\subsection{Conditional expectation and effective filtration convergence}

For a simple function $s$ at level $m$ and a target level $n$, let
$r=\max(m,n)$ and let $\widetilde s$ be its lift to $A_r$.  Define a
rational list on $A_n$ by
\begin{equation}
 (E_ns)_i=
 \begin{cases}
 \displaystyle\frac{\sum_{\pi_{n,r}(j)=i}w_{r,j}\widetilde s_j}{w_{n,i}},
       &w_{n,i}>0,\\[6pt]
 0,&w_{n,i}=0.
 \end{cases}
 \label{eq:m-condexp}
\end{equation}
The branch condition is an exact rational test.  If $w_{n,i}=0$, every
descendant weight is zero, since they are nonnegative and sum to it.
Thus the chosen value zero on that cell has no effect on any $L^1$
observation.  This zero-mass argument is part of the algorithm's proof,
not an implicit division convention.

\begin{theorem}[Conditional expectation on the completion]
\label{thm:m-filtration}
The finite operators \eqref{eq:m-condexp} extend to integrable names as
positive linear contractions.  For all $m,n$,
\begin{equation}
 E_mE_n=E_{\min(m,n)},\qquad I(E_nf)=I(f).
 \label{eq:m-tower}
\end{equation}
If $s_k$ is a regular simple name for $f$, and $s_k$ is stored at level
$r(k)$, then
\begin{equation}
 n\ge r(k+1)\quad\Longrightarrow\quad
 \|E_nf-f\|_1\le\delta_k.
 \label{eq:m-upward-rate}
\end{equation}
The family $(E_nf)_n$ has a uniform-integrability modulus derived directly
from $(s_k)$, without reciprocal lower bounds on nonzero atom masses.
\end{theorem}

\begin{proof}
First work on finite rational lists.  On a positive-mass cell $i$ the
weighted triangle inequality gives
\[
 w_{n,i}|(E_ns)_i-(E_nt)_i|
 \le\sum_{\pi_{n,r}(j)=i}w_{r,j}|s_j-t_j|.
\]
On a zero-mass cell both sides are zero.  Add over all $i$ and regroup
the finite double sum to obtain
\begin{equation}
 \|E_ns-E_nt\|_1\le\|s-t\|_1.
 \label{eq:m-ce-contraction}
\end{equation}
Linearity, positivity, and preservation of the integral follow from the
same cell formula; the zero-mass branch is harmless in each identity.
For $m\le n$, averaging the level-$n$ cell numerators over a level-$m$
cell counts each fine atom once, and gives exactly the level-$m$
numerator.  Division by a positive coarse weight gives $E_mE_n=E_m$;
zero coarse weight again gives zero distance.  Also $E_ms$ is already
constant on the finer level-$n$ cells, so $E_nE_ms=E_ms$ in $L^1$.
This proves the tower property.  The numerator regrouping is precisely
the reusable finite content of \texttt{TowerCommonDen.regroup}; the
normalization and zero-cell argument complete the operator statement.

By \eqref{eq:m-ce-contraction}, the sequence $E_ns_k$ is a regular name
with the original error budget, uniformly in the numerical parameter
$n$.  Lemma~\ref{lem:m-extension} therefore extends the operators and
their identities.  For $n\ge r(k)$ one has $E_ns_k=s_k$ in $L^1$.
Consequently
\[
 \|E_nf-f\|_1
 \le\|E_n(f-s_k)\|_1+\|s_k-f\|_1
 \le2\delta_k.
\]
Replace $k$ by $k+1$ to obtain \eqref{eq:m-upward-rate}.  If a monotone
rate is desired, replace $r(k+1)$ by its maximum over the first $k+2$
approximation levels.  This is a bounded PR operation.

Finally put $K_k=\max(1,\lceil\|s_k\|_\infty\rceil)$.  Each conditional
average of $s_k$ lies in $[-K_k,K_k]$, including the zero-mass branch.
For any bounded $b$ with $|b|\le K$, the scalar inequality
$(|a|-K)^+\le|a-b|$ yields
\[
 T_{K_k}(E_nf)\le\|E_nf-E_ns_k\|_1
                    \le\|f-s_k\|_1\le\delta_k.
\]
Taking bounded maxima of the $K_k$ gives the required monotone cutoff
function.  The contraction proof cancels the cell denominator against
the cell's mass; consequently no smallest positive atom weight enters
this $L^1$ precision budget.
\end{proof}

These operators are genuine conditional expectations on the represented
finite subalgebras.  For every level-$n$ event $B$, finite regrouping gives
$I(\mathbf1_BE_nf)=I(\mathbf1_Bf)$.  An element in the completed
level-$n$ range with these cell integrals is unique in $L^1$: on each
positive cell its constant value is determined by the cell integral,
and zero cells contribute no norm.  This identifies the operator with
the usual conditional expectation in the interval realization.
It does not appeal to a general Radon--Nikodym theorem for an arbitrary
sub-$\sigma$-algebra.

There is a useful distinction between observations.  To request the
conditional value at one specified positive atom to error $\eta$, the
input integral over that atom must be known to error $w_{n,i}\eta$.
That request incurs the expected reciprocal-mass precision loss.  To
request the entire conditional expectation to error $\eta$ in $L^1$,
the contraction \eqref{eq:m-ce-contraction} incurs no such loss.

\paragraph{An unbounded worked example.}
Let
$A_j=[2^{-2j-2},2^{-2j-1})\subset[0,1)$ and let $s_n$ have value $2^j$
on $A_j$ for $j<n$, and zero elsewhere.  These sets are disjoint and
$\mu(A_j)=4^{-j-1}$.  The finite geometric-sum identity gives
\[
 \|s_m-s_n\|_1
   =2^{-n-1}\bigl(1-2^{-(m-n)}\bigr)\quad(m\ge n).
\]
They therefore define an integrable, unbounded element $g$ with
$I(g)=1/2$ and $\|g-s_n\|_1=2^{-n-1}$.  Use the tower whose level $n$
exposes $A_0,\ldots,A_{n-1}$ and a single remaining atom $B_n$.  Its mass
is $(2+4^{-n})/3$.  At level $n$, conditional expectation has the original
values on the exposed annuli and the rational value
\[
 (E_ng)|_{B_n}
   =\frac{3\,2^{-n-1}}{2+4^{-n}}.
\]
Indeed, the unexposed integral is $2^{-n-1}$, as follows from the same
finite geometric identity and its explicit remainder.  On every
unexposed annulus the value of $g$ exceeds this conditional value, while
on the complement of all annuli $g=0$; that complement has mass $2/3$.
Adding these two contributions gives the exact estimate
\[
 \|E_ng-g\|_1=\frac{2^{1-n}}{2+4^{-n}}\le2^{-n}.
\]
The family $(s_n)$ is dominated by $g$.  The extracted cutoff
$U(k)=\max(1,2^{k-1})$ and the rate $R(k)=k$ meet the Vitali hypotheses,
giving the explicit sufficient rate $N(k)=2k+4$.  This example also
illustrates a concrete efficiency gain: the adapted tower has $n+1$
atoms, whereas a uniform dyadic grid resolving the first $n$ annuli has
$4^n$ cells.  The exact-arithmetic program in the archive checks the
finite tower laws, the cutoff budgets and the displayed conditional
values, including independent finite approximations to the last norm.

\begin{corollary}[Representation of a quantitatively convergent martingale]
\label{cor:m-martingale}
Let $M_n$ be uniformly PR-presented integrable elements satisfying
$E_nM_n=M_n$ and $E_mM_n=M_m$ for $m\le n$.  If they satisfy
\eqref{eq:m-ui} and \eqref{eq:m-rho-cauchy}, the name $f$ constructed in
Theorem~\ref{thm:m-vitali} satisfies $M_n=E_nf$ for every $n$.  It is the
unique such integrable element in this completion.
\end{corollary}

\begin{proof}
For fixed $n$ and every $j\ge n$, $E_nM_j=M_n$.  The contraction gives
$\|E_nf-M_n\|_1\le\|f-M_j\|_1$, which tends to zero with the explicit
Vitali schedule.  Hence $E_nf=M_n$.  If also $E_ng=M_n$ for all $n$,
Theorem~\ref{thm:m-filtration} gives $E_nf\to f$ and $E_ng\to g$,
with their respective name-dependent rates.  Comparing both limits to
one sufficiently late common $M_n$ gives $f=g$.
\end{proof}

\begin{proposition}[Why the convergence rate is additional data]
\label{prop:m-martingale-barrier}
There is a uniformly PR family of rational martingales, all bounded by
one, from which no uniform computable $L^1$ Cauchy modulus can be
extracted.
\end{proposition}

\begin{proof}
At level $n$ of the fair-bit probability tower expose bits
$\omega_1,\ldots,\omega_n$.  Delay each machine by one initial step, so
its first possible halting step is at least one.  For a machine index $e$, simulate its first
$n$ steps.  Put $M^e_n=0$ if it has not halted.  If its first halting step
is $t\le n$, put $M^e_n(\omega)=(-1)^{\omega_t}$.  This is a rational
table produced by a bounded simulation and bounded search, uniformly PR
in $e,n$ and the atom index.  If the machine has halted, the martingale
stays unchanged.  If it first halts at the next step, its new value is the
new independent fair sign, with conditional average zero.  In every
case $E_nM^e_{n+1}=M^e_n$.  The absolute bound is one, so the excess-tail
modulus is constantly one with tail zero.

Suppose a computable procedure supplied $N(e)$ such that
$\|M^e_m-M^e_n\|_1\le1/2$ for all $m,n\ge N(e)$.  Simulate $e$ through
$N(e)$ steps.  If it has not halted, it can never halt later: at a later
first halting step $t$, $M^e_{N(e)}=0$ and $|M^e_t|=1$, whose $L^1$
distance is one.  Thus this procedure would decide the halting problem.
The contradiction proves the claim.  The same argument rules out a
uniform $\rho$ Cauchy modulus, since the jump also has truncated distance
one.  It leaves intact the explicit modulus in
\eqref{eq:m-upward-rate}, whose input is already a regular integrable
name of the terminal object.
\end{proof}

\subsection{Exact internalization obligations}

The mathematical induction used above is finite: regrouping sums,
normalizing positive rational denominators, looping over atom tables,
taking bounded maxima, and proving the power-of-two budget.  Each has a
quantifier-free numerical invariant.  For example, the clipping endpoint
after multiplying by positive common denominators is the universal
closure of
\[
 K\ge1\ \Longrightarrow\
 \sum_{i<a(r)}w_i|x_i-y_i|
 \le\sum_{i<a(r)}w_ip_K(x_i)
    +\sum_{i<a(r)}w_ip_K(y_i)
    +2K\sum_{i<a(r)}w_i\min(1,|x_i-y_i|).
\]
Every symbol in this formula abbreviates a fixed PR term on numerical
codes.  The corresponding output must ultimately have type
\texttt{Prf PRAAx []} applied to the encoded formula.  A Lean proof about
$\Q$-valued functions is a useful independently checked semantic kernel,
but is not itself that derivation.

The remaining formalization path is specific.  First instantiate the
repository's rational normalization and finite sum interfaces for the
tower codes, including the zero-mass branch.  Prove the scalar clipping
and excess-Lipschitz inequalities by rational case splits; sum them using
the existing order lemmas.  Reify $U,R,\ell,q,N$ and the simple output
program \eqref{eq:m-output-name}.  For each requested finite observation,
substitute these terms into \eqref{eq:m-pair-cauchy} and the input
regularity formulas.  This produces output regularity, convergence, and
the inherited-tail observations by ordinary logical rules.  Finally
instantiate \texttt{TowerCommonDen.regroup}, prove the normalized
contraction, and obtain \eqref{eq:m-upward-rate} from the stated level
index.  None of these steps requires an induction hypothesis asserting
convergence at every precision.  Precision is a free numerical
parameter; the finite loop and all finer demands are explicit terms.

The interval model supplies semantic interpretation of the resulting
arithmetic laws.  The laws apply to the chosen completion and its
represented subalgebras.  Finite measures follow by an explicitly
positive normalization; extensions to an infinite ambient measure also
need a spatial-tail or exhaustion contract.  General pointwise
convergence, arbitrary sub-$\sigma$-algebras, and arbitrary countable
unions do not silently provide any of those effective data.
